The result that a solution to the LR ESLP provides a lower bound on the ESLP
objective function and that cost associated with any feasible $(x, y)$ provides
an associated upper bound motivates an approach to obtaining an approximate solution to the 
ESLP through iteratively obtaining solutions that tighten lower and upper bounds. Such heuristics
have been suggested as efficient approaches to solving large scale instances 
of the CFLP \cite{cornuejols1991comparison}.

A basic outline of the heuristic proposed is now described. The approach is to iteratively improve
the lower bound estimation obtained via the LR ESLP by generating a sequence of 
$\lambda$ parameters in order to achieve the highest lower bound possible. This
is achieved through implementing a subgradient optimisation algorithm that updates those $\lambda$
parameters corresponding to violations of the original ESLP. A subgradient method is chosen due to
its simplicity and widespread use in LR frameworks, and the specific subgradient algorithm chosen
is the classical subgradient method, also implemented in \cite{wu2006capacitated}. In future, this
is an aspect that could be modified to investigate possible performance improvements. Other 
papers also find that subgradient based
heuristics can also achieve good-quality solutions \cite{cornuejols1991comparison}.
Each iteration of this method also presents an opportunity to obtain an improved upper bound, as from each LR ESLP solution
a feasibible solution to the ESLP can be obtained. If this corresponds to
a lower cost than those found previously, then an improved upper bound of the ESLP is also found.
The relative gap between these upper and lower bounds is used as both a stopping criterion for this 
heuristic and a measure of error.

\begin{equation}\label{eq:gap}
    \text{gap} = \frac{UB - LB}{UB}
\end{equation}


\subsubsection{The lower bound estimation}

Each supply site $s$ at time $t$ in the LR ESLP may violate the capacity constraint, resulting in a 
subgradient component associated with that constraint. Specifically, let

\begin{equation} g_{s,t} = \sum_{d \in D} x_{s,d,t} - u_s y_s \end{equation}

denote the violation of the capacity constraint at $(s,t)$. Note that $g_{s,t}$ is derived from 
the current LR ESLP solution $(x,y)$ and indicates whether the associated capacity constraint is 
respected ($g_{s,t} \leq 0$) or violated ($g_{s,t} > 0$).

We initialise all Lagrange multipliers $\lambda_{s,t}$ at $0$, corresponsing to no initial penalisation.
This was chosen as empirically resulted in good convergence.
For each violation i.e. case where $g_{s,t} \geq 0$, we update
the lagrange multiplier $\lambda_{s,t}$ according to the following subgradient rule. First, we define step-size as follows:

\begin{equation}\label{eq:subgradient}
    \alpha_k = \frac{UB_k - LB_k}{\sum_{s,t} g_{s,t}^2} 
\end{equation}

This step size scales inversely with the magnitude of the subgradient,
ensuring more stable updates as the solution improves.

For each violation (i.e., whenever $g_{s,t} > 0$), we then update the corresponding Lagrange multiplier:

\begin{equation}
    \lambda_{s,t}^{k+1} = \max\bigl(0, \lambda_{s,t}^{k} + \alpha_k g_{s,t}\bigr).
\end{equation}

The $\max$ function ensures the multipliers remain non-negative, and 
therefore ensure that we penalise constraint violations. By iteratively updating the multipliers based on the subgradients
the LR ESLP solutions increasingly respect the original capacity constraints.

\subsubsection{The upper bound estimation}

If the solution to the LR ESLP $(x_{LR}, y_{LR})$ is feasible, we can use it to obtain an upper bound
by computing the corresponding cost implied by putting the solution into the ESLP's objective function 
\ref{eq:ESLP_obj}. Typically, however, the LR ESLP does not find a feasible solution though, and so we 
perform a feasibility adjustment to $(x_{LR}, y_{LR})$ to obtain a feasible solution $(x, y)$, 
which is then used to obtain an upper bound on the ESLP. \par

Each violation can be decomposed into a subproblem focussed on a specific supply site $s$ at time $t$
that provides an excees capacity defined as:

\begin{equation}
    \text{Excess}_{s,t} = max(\sum_{d \in D}^{} x_{LR, s,d,t} - u_sy_{LR s}, 0).
\end{equation}


To adjust, a greedy algorithm is implemented which redistributes the excess to demand sites in a cost-effective
manner. For each subproblem, the demand sites $d \in D$ are arranged in ascending order based on 
their allocation cost $c_{s, d, t}$, and each site then allocates remaining supply until the demand at demand site $d$  
is either matched or supply site $s$ has no remaining supply to provide, in which case we move to the 
next supply site. The algorithm is outlined in \ref{alg:greedy_feasibility}

\begin{algorithm}[H]
    \SetAlgoLined
    \KwIn{Lagrangian Relaxed Solution $(x_{LR}, y_{LR})$, Demand $r_{d,t}$ for all $d \in D$, $t \in T$, Capacity $u_s$ for all $s \in S$, Big-M constant $M$}
    \KwOut{Feasible Solution $(x^*, y^*)$}
    
    \tcp{Initialize Feasible Solution}
    Initialize $(x^*, y^*) \leftarrow (x_{LR}, y_{LR})$\;
    
    \ForEach{$s \in S$}{
        \ForEach{$t \in T$}{
            \tcp{Calculate Excess Supply}
            $\text{Excess}_{s,t} \leftarrow \sum_{d \in D} x_{LR,s,d,t} - u_s y_{LR,s,t}$\;
            
            \If{$\text{Excess}_{s,t} > 0$}{
                \tcp{Sort Demand Sites by Allocation Cost}
                $sorted\_demand \leftarrow$ sort $D$ by $c_{s,d,t}$ ascending\;
                
                \tcp{Allocate Excess Supply}
                \ForEach{$d_i \in sorted\_demand$}{
                    $\text{Available\_Demand} \leftarrow r_{d_i,t} - \sum_{s' \in S} x^*_{s',d_i,t}$\;
                    \If{$\text{Available\_Demand} > 0$}{
                        $\text{Allocation} \leftarrow \min(\text{Excess}_{s,t}, \text{Available\_Demand})$\;
                        $x^*_{s,d_i,t} \leftarrow x^*_{s,d_i,t} - \text{Allocation}$\;
                        $x^*_{s',d_i,t} \leftarrow x^*_{s',d_i,t} + \text{Allocation}$\;
                        $\text{Excess}_{s,t} \leftarrow \text{Excess}_{s,t} - \text{Allocation}$\;
                        
                        \If{$\text{Excess}_{s,t} \leq 0$}{
                            \textbf{break}\;
                        }
                    }
                }
            }
        }
    }
    
    \caption{Greedy Feasibility Adjustment for ESLP}\label{alg:greedy_feasibility}
\end{algorithm}

\subsubsection{Algorithm for LR Heuristic}

Combining the results and steps outlined above, we present the LR Heurisic implemented which
will be assessed in this paper, algorithm \ref{alg:LRHeuristic}.

\begin{algorithm}[H]
    \SetAlgoLined
    \SetKwInOut{Input}{Input}
    \SetKwInOut{Output}{Output}
    
    \Input{Initial values for Lagrange multipliers $\lambda_{s,t} \geq 0$, convergence threshold $\epsilon$, maximum iterations $maxIter$}
    \Output{Final lower bound $LB$, upper bound $UB$, and feasible solution $(x^*, y^*)$}
    
    Initialize $LB \leftarrow -\infty$, $UB \leftarrow \infty$, $k \leftarrow 0$\;
    
    \While{$k \le maxIter$ \textbf{or} $\frac{UB - LB}{UB} \geq \epsilon$}{
        Solve LR ESLP \eqref{eq:LR_ESLP} with current $\lambda_{s,t}$ to obtain solution $(x_{LR}^k, y_{LR}^k)$\;
        Set $LB \leftarrow \max(LB, \mathcal{L}(x_{LR}^k, y_{LR}^k))$\;
        \If{$(x_{LR}^k, y_{LR}^k)$ is not feasible}{
            Obtain feasible $(x^k, y^k)$ using greedy algorithm \ref{alg:greedy_feasibility}\;
        }
        \Else{
            Set $(x^k, y^k) \leftarrow (x_{LR}^k, y_{LR}^k)$\;
        }
        Set $UB \leftarrow \min(UB, \sum_{s \in S} f_s y^k_s + \sum_{s \in S} \sum_{d \in D} \sum_{t \in T} c_{s,d,t} x^k_{s,d,t})$\;
        \If{$\frac{UB - LB}{UB} \leq \epsilon$}{
            Stop.}
        \Else{$\lambda_{s,t} \leftarrow \max(0, \lambda_{s,t} + alpha_k g_{s,t}) \forall s \in S,\forall k \in K$}
        Increment $k \leftarrow k + 1$\;
    }
    \caption{Lagrangian Relaxation Heuristic for Energy System Location Problem}\label{alg:LRHeuristic}
\end{algorithm}


It is highlighted that if the LR ESLP solution is feasible without modification, then we have found the
optimal solution to the ESLP as this corresponds to $LB = UB$. 
The above algorithm terminates in this case instantly.

\newpage
